% IREF working appendix fragment; all numbers copied from frozen V004 outputs.
% Requires booktabs. No replacement for the archived V003/V004 numerical files.
\section{Calibration of the exploratory dimension-comparison intervals}
\label{app:inference-calibration}

Experiment P1-G4-V004 evaluates the previously fixed simultaneous-interval
procedure without fitting additional networks or examining additional empirical
periods. Six candidate labels and five model variants are represented by pricing
moment vectors with known population means. After a fixed whitening transform,
their population distance is $d_{ks}=\sqrt{12}\|\mu_{ks}\|_2$. Deterministic
rotations give the five variants the same population distance within each
candidate. Stationary Gaussian AR(1) innovations combine common, candidate, and
variant components with variance shares 0.5, 0.3, and 0.2. Total innovation energy
is held constant across moment ranks. This is a moment-level calibration,
conditional on fixed models and weights, rather than a structural return/SDF
model or an estimate of power for the empirical panel.

For each simulated sample, we reproduce the original procedure: common
12-month circular time blocks, five jointly resampled variant indices, the
average of variant-specific sample distances, and all 15 pairwise differences.
For difference $j$, let $\widehat\Delta_j$ be its point estimate,
$\widehat\Delta_j^{*(b)}$ its bootstrap value, and $s_j^*$ the bootstrap standard
deviation. The common critical value is the 0.95 bootstrap quantile of
$\max_j |(\widehat\Delta_j^{*(b)}-\widehat\Delta_j)/s_j^*|$; the interval is
$\widehat\Delta_j\pm c_{.95}s_j^*$. This is a single-step maximum standardized
bootstrap-deviation construction. It does not re-estimate the standardizer
inside each draw and does not implement a stepdown procedure.

The frozen design contains 30 scenarios, each with 400 independently seeded
outer replications and 1,999 inner draws. The primary scenarios use 74 moment
coordinates, 120 months, baseline distance five, AR coefficients zero or 0.5,
and population gaps of zero, 0.5, one, or two between the first five candidates
and the sixth. Prespecified sensitivities extend the sample to 240 or 600 months,
reduce the moment rank to four, or set all distances to zero. Candidate labels
do not represent a known structural factor count. The design contains no
post-calibration adjustment to the block length, candidate grid, or interval
construction.

\begin{widetable}[!tbp]
\centering
\caption{Equal-distance calibration: simultaneous coverage and familywise rejection}
\label{tab:inference-calibration}
\small
\begin{tabular}{rrrrrr}
\toprule
Months & Rank & Distance & AR coefficient & Coverage (\%) & FWER (\%) \\
\midrule
120 & 74 & 5 & 0.0 & 83.50 & 16.50 \\
120 & 74 & 5 & 0.5 & 87.00 & 13.00 \\
240 & 74 & 5 & 0.0 & 91.50 & 8.50 \\
240 & 74 & 5 & 0.5 & 91.75 & 8.25 \\
600 & 74 & 5 & 0.0 & 93.25 & 6.75 \\
600 & 74 & 5 & 0.5 & 93.00 & 7.00 \\
120 & 4 & 5 & 0.0 & 86.25 & 13.75 \\
120 & 4 & 5 & 0.5 & 90.50 & 9.50 \\
120 & 74 & 0 & 0.0 & 98.75 & 1.25 \\
120 & 74 & 0 & 0.5 & 98.50 & 1.50 \\
\bottomrule
\end{tabular}
\par\smallskip
\begin{minipage}{0.96\linewidth}\footnotesize
Each row contains 400 outer replications. Coverage requires all 15 intervals
to include their population differences. FWER is the probability of excluding
zero in at least one comparison; since all differences in these rows equal
zero, it equals one minus simultaneous coverage. Nominal values are 95\%
coverage and 5\% FWER. Wilson Monte Carlo 95\% intervals for coverage in the
first two rows are [79.55,86.82]\% and [83.35,89.95]\%; the corresponding FWER
intervals are [13.18,20.45]\% and [10.05,16.65]\%. These probabilities describe
the specified simulations, not empirical error probabilities for U.S. equities.
All 30 scenarios, including unequal-distance detection results, are retained
in the experiment archive.
\end{minipage}
\end{widetable}

The primary short-sample equal-distance scenarios exhibit undercoverage and
excess familywise rejection. Longer samples improve coverage in those scenarios
but do not establish exact calibration. The two zero-distance stress scenarios
are conservative, so the primary failures cannot be attributed simply to the
norm's nondifferentiability at zero. The experiment does not isolate the
contributions of block length, nonlinear transformation, dependence, and
variant resampling. In particular, the fixed rotated variants do not reproduce
the distribution of repeatedly trained networks. These limits restrict the
interpretation of the calibration without removing its adverse result.

This exercise does not calibrate the separate predictive-$R^2$ intervals,
the pointwise difference from FF5 plus momentum, or the earlier paper's other
procedures. Nor does failure of this particular construction establish a new
economic mechanism. Its consequence is narrower: the exploratory simultaneous
intervals are not used to support equivalence, a confidence set for a unique
dimension, or structural nonidentification. Point estimates and all original
outputs remain reported without retuning the empirical models.
